% Pandoc template for IEEE conference paper sections
\section{Conclusion}

We presented three contributions to pairing-friendly elliptic curve
construction:

\begin{enumerate}
\def\labelenumi{\arabic{enumi}.}
\item
  \textbf{NTT-friendly Cocks-Pinch.} By constraining the cyclotomic seed
  \(T \equiv 0 \pmod{2^{12}}\), we guarantee two-adicity \(\geq 36\) in
  the scalar field deterministically, while an extended lift search
  achieves the same for the base field. This exceeds the
  industry-standard BLS12-381 (two-adicity 32) at a modest increase in
  generation time (30--60 s vs.~2--10 s).
\item
  \textbf{Readability-aware prime selection.} Our multi-criteria scoring
  system produces primes that are more amenable to human inspection
  without any security degradation, addressing the auditability concern
  in cryptographic parameter generation.
\item
  \textbf{Curve1024.} We instantiate these techniques into a concrete,
  publicly verifiable curve: \(E: y^2 = x^3 + 41\) over a 1024-bit prime
  field with embedding degree 18, providing 256-bit ECDLP security and
  NTT support up to length \(2^{36}\) in both fields. The complete
  parameter set (field primes, generator point, CM data) is published
  and validated by a Rust implementation with 36 passing tests.
\end{enumerate}

The system is suitable for offline signing applications requiring
long-term security guarantees and for future pairing-based protocols
(BLS aggregate signatures, zk-SNARKs) that benefit from the NTT-friendly
field structure.

\textbf{Future work} includes: (1) projective/Jacobian coordinates to
eliminate per-operation inversions, (2) constant-time Montgomery ladder
for side-channel resistance, (3) full optimal Ate pairing implementation
over \(\mathbb{F}_{p^{18}}\), and (4) extension to other embedding
degrees (\(k = 12, 24\)).

\subsection{Availability}\label{availability}

The implementation and all curve parameters are available as open-source
software under the MIT license \cite{LumenMath}.
